\documentclass[11pt]{scrartcl}

\usepackage{amsmath,amsthm,amssymb,mathtools}
\usepackage{hyperref}
\usepackage{lastpage}
\usepackage{multicol}

\theoremstyle{plain} 
\newtheorem{theorem}{Theorem}[section]
\newtheorem{lemma}[theorem]{Lemma}
\newtheorem*{question}{Question}
\theoremstyle{definition}
\newtheorem*{definition}{Definition} 
\theoremstyle{remark}
\newtheorem*{remark}{Remark}

\usepackage[none]{hyphenat}
\setlength{\parindent}{0pt}
% setlength{\emergencystretch}{3em}

\renewcommand\labelitemi{$\vcenter{\hbox{\tiny$\bullet$}}$} 

\usepackage[margin=1in]{geometry}
\usepackage{fancyhdr}
\pagestyle{fancy}
\lhead{ESC103: Engineering Mathematics and Computation} 
\rhead{\thepage}
\cfoot{}

\title{ESC103: Engineering Mathematics and Computation \\[0.1in] Exam Review}
\author{Brendan Tang}
\date{Fall 2026}

\begin{document}

\maketitle
\thispagestyle{empty}

\tableofcontents

\newpage
\section{Vectors}
\subsection{Vector Operations}

Here are some basic vector operations in $\mathbb{R}^{2}$. 

\begin{itemize}
    \item Vector Addition: $\begin{bmatrix}v_{1} \\ v_{2}\end{bmatrix}+\begin{bmatrix} w_{1} \\ w_{2}
    \end{bmatrix} = \begin{bmatrix}v_{1}+w_{1} \\ v_{2}+w_{2}\end{bmatrix}$
    \item Scalar Multiplication: $c\begin{bmatrix}v_{1} \\ v_{2}\end{bmatrix}=\begin{bmatrix}cv_{1} \\ cv_{2}
    \end{bmatrix}$
\end{itemize}

Both are commutative, associative, and scalar multiplication is distributive over vector addition. \\[-0.5em]

There are two ways to multiply vectors in $\mathbb{R}^{3}$. The dot product takes two vectors and returns a scalar. 
Here are two definitions for the dot product.

\begin{definition}
    Let $\mathbf{v}=\begin{bmatrix}v_{1} \\ v_{2} \\ v_{3}\end{bmatrix}$ and $\mathbf{w}=\begin{bmatrix}w_{1} \\ 
    w_{2} \\ w_{3}\end{bmatrix}$. The algebraic definition of the dot product is
    \[ \mathbf{v}\cdot \mathbf{w}=v_{1}w_{1}+v_{2}w_{2}+v_{3}w_{3}. \]
    Let $\theta$ be the angle between $\mathbf{v}$ and $\mathbf{w}$. The geometric definition of the dot product is 
    \[ \mathbf{v}\cdot \mathbf{w}=\lVert \mathbf{v} \rVert \lVert \mathbf{w} \rVert \cos \theta. \]
\end{definition}

The dot product is commutative, distributive over vector addition, and not associative. \\[-0.5em]

The cross product takes two vectors and returns a vector that is orthogonal to both. Here are two definitions for 
the cross product.

\begin{definition}
    Let $\mathbf{v}=\begin{bmatrix}v_{1} \\ v_{2} \\ v_{3}\end{bmatrix}$ and $\mathbf{w}=\begin{bmatrix}w_{1} \\ 
    w_{2} \\ w_{3}\end{bmatrix}$. The algebraic definition of the cross product is
    \[ \mathbf{v}\times \mathbf{w}=\begin{bmatrix}v_{2}w_{3}-v_{3}w_{2} \\ v_{3}w_{1}-v_{1}w_{3} \\ 
    v_{1}w_{2}-v_{2}w_{1} \end{bmatrix}. \]
    Let $\theta$ be the angle between $\mathbf{v}$ and $\mathbf{w}$. The geometric definition of the cross product is 
    \[ \lVert \mathbf{v}\times \mathbf{w} \rVert=\lVert \mathbf{v} \rVert \lVert \mathbf{w} \rVert \sin \theta. \]
\end{definition}

The cross product is anti-commutative, not distributive, and not associative.

\subsubsection{Vector Projection}

The normalization of a vector, $\mathbf{v}$ is
\[ \hat{\mathbf{v}}=\frac{\mathbf{v}}{\lVert \mathbf{v} \rVert }. \]
The scalar projection of $\mathbf{v}$ onto $\mathbf{w}$ is 
\[ \text{comp}_{\mathbf{w}}\mathbf{v} = \frac{{\mathbf{v}\cdot \mathbf{w}}}{\lVert \mathbf{w} \rVert }\ 
\text{(algebraic)}\quad \text{comp}_{\mathbf{w}}\mathbf{v} = \lVert \mathbf{v} \rVert \cos \theta\ \text{(geometric)}.  \]
Combining both, the formula for the vector projection of $\mathbf{v}$ onto $\mathbf{w}$ is
\[ \text{proj}_{\mathbf{w}}\mathbf{v} = \frac{{\mathbf{v}\cdot \mathbf{w}}}{\lVert \mathbf{w}^{2} \rVert }\mathbf{w}. \]

\subsection{Vector Properties}

Here are some special types of vectors in $\mathbb{R}^{3}$.

\begin{itemize}
    \item Zero Vector: A vector with length $0$, denoted by $\mathbf{0}=\begin{bmatrix}0 \\ 0 \\ 0\end{bmatrix}$
    \item Unit Vector: A vector with length $1$, denoted by $\hat{\mathbf{v}}$
    \item Standard Basis Vectors: Special unit vectors denoted by $\hat{\mathbf{i}}=\begin{bmatrix}1 \\ 0 \\ 0
    \end{bmatrix}$, $\hat{\mathbf{j}}=\begin{bmatrix}0 \\ 1 \\ 0\end{bmatrix}$, and $\hat{\mathbf{k}}=\begin{bmatrix}
    0 \\ 0 \\ 1\end{bmatrix}$
\end{itemize}

The length of a vector $\mathbf{v}$ in $\mathbb{R}^{3}$ is 
\[ \lVert \mathbf{v} \rVert  =\sqrt{ v_{1}^{2}+v_{2}^{2}+v_{3}^{2} }. \]

Two vectors $\mathbf{v}$ and $\mathbf{w}$ are orthogonal if
\[ \mathbf{v}\cdot \mathbf{w}=0. \]

Two vectors $\mathbf{v}$ and $\mathbf{w}$ are collinear if
\[ \mathbf{v}\times \mathbf{w}=\mathbf{0}. \]

\subsection{Plane Geometry}

A line in $\mathbb{R}^{3}$ can be represented by the vector equation
\[ \begin{bmatrix}x \\ y \\z\end{bmatrix}=\begin{bmatrix}x_{0} \\ y_{0} \\z_{0}\end{bmatrix}+a\mathbf{d}\quad a\in 
\mathbb{R}, \]
where $\mathbf{d}$ is a direction vector and $(x_{0},y_{0},z_{0})$ is a point on the line. \\[-0.5em]

Similarly, a plane in $\mathbb{R}^{3}$ can be represented by the vector equation
\[ \begin{bmatrix}x \\ y \\z\end{bmatrix}=\begin{bmatrix}x_{0} \\ y_{0} \\z_{0}\end{bmatrix}+a\mathbf{d}_{1}+
b\mathbf{d}_{2}\quad a,\in \mathbb{R}, \]
where $\mathbf{d}_{1}$, $\mathbf{d}_{2}$ are non-collinear direction vectors and $(x_{0},y_{0},z_{0})$ is a point 
on the plane. \\[-0.5em]

A plane in $\mathbb{R}^{3}$ can also be represented by the scalar equation
\[ ax+by+cz=d, \]
where $\mathbf{n}=\begin{bmatrix}a \\ b \\ c\end{bmatrix}$ is a vector normal to the plane. (why?).

\subsubsection{Interlude (Row and Column Pictures)}

The row picture of a system of linear equations treats each equation as a geometric line or plane, and the solution 
is the point where they all intersect. \\[-0.5em]

The column picture of a system of linear equations views the system as a linear combination of column vectors, and 
to solve you need to find the correct scalar multiple of each vector. 

\newpage
\section{Matrices}
\subsection{Matrix Operations}

Here are some basic matrix operations.

\begin{itemize}
    \item Matrix Addition: $\begin{bmatrix}a_{11} & a_{12} \\ a_{21} & a_{22}\end{bmatrix}$ + $\begin{bmatrix}
    b_{11} & b_{12} \\ b_{21} & b_{22}\end{bmatrix}$ = $\begin{bmatrix}a_{11} + b_{11} & a_{12} + b_{12} \\ 
    a_{21} + b_{21} & a_{22} + b_{22}\end{bmatrix}$
    \item Scalar Multiplication: $c\begin{bmatrix}a_{11} & a_{12} \\ a_{21} & a_{22}\end{bmatrix}=\begin{bmatrix}
    ca_{11} & ca_{12} \\ ca_{21} & ca_{22}\end{bmatrix}$
\end{itemize}

Both are commutative, associative, and scalar multiplication is distributive over matrix addition. \\[-0.5em]

If $\mathbf{A}$ is an $m\times r$ matrix and $\mathbf{B}$ is an $r\times n$ matrix, we can perform matrix 
multiplication $\mathbf{A}\mathbf{B}$ to get an $m\times n$ matrix.

\begin{definition}
    Let $\mathbf{A}=\begin{bmatrix}\text{---}\ \mathbf{a}_{1}\ \text{---} \\ \vdots \\ \text{---}\ \mathbf{a}_{n}\ \text{---}
    \end{bmatrix}$ and $\mathbf{B}=\begin{bmatrix}| & & | \\ \mathbf{b}_{1} & \cdots & \mathbf{b}_{2} \\ | & & |\end{bmatrix}$. 
    Then the row-column definition is
    \[ \mathbf{A}\mathbf{B}=\begin{bmatrix}\mathbf{a}_{1}\cdot \mathbf{b}_{1} & \cdots & \mathbf{a}_{1}\cdot \mathbf{b}_{n} \\ 
    \vdots & & \vdots \\ \mathbf{a}_{n}\cdot \mathbf{b}_{1} & \cdots & \mathbf{a}_{n}\cdot \mathbf{b}_{n}\end{bmatrix},\]
    where the value in the $i$th row and $j$th column of the resultant matrix is equal to $\mathbf{a}_{i}\cdot 
    \mathbf{b}_{j}$. \\[-0.5em]

    Let $\mathbf{A}=\begin{bmatrix}| & & | \\ \mathbf{a}_{1} & \cdots & \mathbf{a}_{n} \\ | & & |\end{bmatrix}$ and $\mathbf{B}=
    \begin{bmatrix}\text{---}\ \mathbf{b}_{1}\ \text{---} \\ \vdots \\ \text{---}\ \mathbf{b}_{n}\ \text{---}\end{bmatrix}$.
    Then the column-scaling definition is
    \[ \mathbf{A}\mathbf{B}=\begin{bmatrix}b_{11}\begin{bmatrix}| \\ \mathbf{a}_{1} \\ |\end{bmatrix}+\dots +b_{n1}
    \begin{bmatrix}| \\ \mathbf{a}_{1} \\ |\end{bmatrix} &\cdots &b_{1n}\begin{bmatrix}| \\ \mathbf{a}_{n} \\ |\end{bmatrix}
    +\dots + b_{nn}\begin{bmatrix}| \\ \mathbf{a}_{n} \\ |\end{bmatrix}\end{bmatrix}, \]
    where the $j$th column of the resultant matrix is a linear combination of $\mathbf{a}_{n}$ using $b_{nj}$ as 
    coefficients.
\end{definition} % This entire definition was a nightmare to type

Matrix multiplication is associative, distributive over vector addition, and not commutative. 

\subsubsection{Geometric Solutions to $\mathbf{A}\mathbf{x}=\mathbf{b}$}

We can use row and column pictures to multiply a matrix by a vector. Consider the equation
\[ \begin{bmatrix}a_{11} & a_{12} & a_{13} \\ a_{21} & a_{22} & a_{23} \\ a_{31} & a_{32} & a_{33}\end{bmatrix}
\begin{bmatrix}x_{1} \\ x_{2} \\ x_{3}\end{bmatrix}=\begin{bmatrix}b_{1} \\ b_{2} \\ b_{3}\end{bmatrix}. \]

\begin{itemize}
    \item Row Picture: We get three planes in $\mathbb{R}^{3}$,
    \[ \begin{aligned} a_{11}x_{1}+a_{12}x_{2}+a_{13}x_{3}&=b_{1} \\ a_{21}x_{1}+a_{22}x_{2}+a_{23}x_{3}&=b_{2} \\ 
    a_{31}x_{1}+x_{32}x_{2}+a_{33}x_{3}&=b_{3}. \end{aligned} \]
    \item Column Picture: We get a linear combination of three vectors in $\mathbb{R}^{3}$,
    \[ x_{1}\begin{bmatrix}a_{11} \\ a_{21} \\ a_{31}\end{bmatrix} + x_{2}\begin{bmatrix}a_{12} \\ a_{22} \\ a_{32}
    \end{bmatrix} + x_{3}\begin{bmatrix}a_{13} \\ a_{23} \\ a_{33}\end{bmatrix} = \begin{bmatrix}b_{1} \\ b_{2} \\ 
    b_{3}\end{bmatrix}. \]
\end{itemize}

\begin{question}
    For each interpretation, what does it mean geometrically for there to be exactly one solution 
    $(x_{1},x_{2},x_{3})$? Infinitely many solutions? No solutions?
\end{question}

\subsubsection{Interlude (Important Definitions)}

Here are some important definitions in linear algebra.

\begin{itemize}
    \item Linear Independence: $\{\mathbf{a}_{n}\}$ are independent if the only solution to $\mathbf{A}
    \mathbf{x}=\mathbf{0}$ is $\{x_{n}\}=0$
    \item Basis: A set of linearly independent vectors that span the vector space
    \item Span: The span of a basis is the set of all linear combinations of the basis % Circular???
    \item Dimension: The dimension of a space is the number of vectors in any basis for that space
    \item Column Space: The column space of a matrix is the span of its column vectors
    \item Rank: The rank of a matrix is the dimension of its column space
\end{itemize}

\subsection{Column-Row Factorization}

Column-row factorization decomposes any $m\times n$ matrix $\mathbf{A}$ of rank $r$ into the product of an $m\times 
r$ matrix $\mathbf{C}$ and an $r\times n$ matrix $\mathbf{R}$. 

\end{document}